% Source: physical PDF pp. 416--420 (printed pp. 393--397), from the
% Section 32.2 heading through the final sentence immediately before
% Section 32.3 on physical p. 420.
% Inventory: Eq. (32.2.1), one unnumbered three-line display group,
% linked citations [7] and [4], and no footnotes, figures, tables, or
% typographic dividers.

\subsection{Massless Multiplets}

We will now consider how the supersymmetry algebras constructed in the
previous section may be used to construct supermultiplets of massless
particle states in \(d\geq4\) spacetime dimensions. The momentum operator
\(P^\mu\) commutes with all fermionic symmetry generators, so we can work
in the one-particle subspace of Hilbert space where \(P^\mu\) has a definite
lightlike eigenvector \(p^\mu\), which may be taken in a direction with
\(p^1=p^0\) and all other spatial components of \(p^\mu\) zero. Just as in
the case of four spacetime dimensions, discussed in Section 2.5, these
one-particle states are classified according to the finite-dimensional
representations they provide of the little group, the subgroup of the
homogeneous Lorentz group that leaves \(p^\mu\) invariant. The little group
contains combined boosts along directions perpendicular to \(\mathbf p\)
and rotations in planes in which \(\mathbf p\) lies, such as the
transformations (2.5.6) for four spacetime dimensions, but these form an
invariant Abelian subgroup and therefore in a finite-dimensional
representation must be represented by the unit operator. With this
subgroup omitted, the reduced little group in \(d\) dimensions is
\(O(d-2)\), consisting of rotations in planes orthogonal to \(\mathbf p\).
We therefore classify massless particle states according to the
representations they provide of \(O(d-2)\) and of the automorphism groups
described at the end of the previous section.

These representations are more complicated than those we have dealt with
in four spacetime dimensions, where the reduced little group is \(O(2)\)
and the representations are one-dimensional, characterized by a single
number, the helicity. Nevertheless it is useful to label the
representations of the reduced little group \(O(d-2)\) with a `spin,'
defined as the \emph{maximum} absolute value of the eigenvalue of any
generator \(J_{ij}\) in the representation.

It is widely believed that there are no consistent quantum field theories
involving massless particles with spin greater than 2. It is
known~\hyperref[ch32-ref-7]{[7]} that \emph{soft} massless particles with
spin \(j>1/2\) can only interact with conserved currents carrying spin
\(j\). For \(j=1\) these are the currents of ordinary conserved scalars,
like electric charge; for \(j=3/2\) they are the one or more supercurrents
associated with supersymmetry; for \(j=2\) there is a single current, the
energy-momentum tensor; but for \(j>5/2\) there is no conserved current
with which a soft massless particle could interact. We may derive
stringent limits on the dimensions in which supersymmetry is possible by
adopting a prohibition against there being more than one type of massless
particle of spin 2, or any massless particles whatever of spin greater
than 2.

Let us return for a moment to the classification of operators used in
Section 32.1, according to a weight equal to the value of the \(O(d)\)
generator \(J_{d1}\) that they destroy. (Recall that
\(J_{01}=iJ_{d1}\).) The fermionic supersymmetry generators have weights
\(1/2\) and \(-1/2\), so the anticommutator of any of these generators
with its Hermitian adjoint can only have weight \(+1\) or \(-1\),
respectively, and therefore must be respectively proportional to the
operator \(P^0+P^1\) or \(P^0-P^1\). But we are working in a subspace of
Hilbert space in which the operator \(P^0-P^1\) vanishes, so in this
subspace the fermionic supersymmetry generators of weight \(-1/2\) all
vanish. To classify the one-particle states we therefore have available
just half of the supersymmetry generators, the \(2^{n-1}\) generators with
weight \(\sigma_{d1}=+1/2\).

We can further divide the remaining supersymmetry generators into two
classes, those in which \(\sigma_{23}=1/2\) or
\(\sigma_{23}=-1/2\) as well as \(\sigma_{d1}=+1/2\). Since the operator
\(P^0+P^1\) has \(\sigma_{23}=0\), the fermionic supersymmetry generators
of each class anticommute with each other, though not necessarily with
their adjoints or with generators of the other class.

Now consider a representation of the little group \(O(d-2)\) with spin
\(j\), and consider any state \(\lvert\lambda\rangle\) that is an
eigenstate of \(J_{23}\) with eigenvalue \(\lambda>0\) and is annihilated
by all supersymmetry generators with \(\sigma_{23}=-1/2\). (Any state that
has the maximum eigenvalue \(j\) for \(J_{23}\) is of this type, but in
general there may be other such states.) We may form states with
\(J_{23}=\lambda-k/2\) by acting on \(\lvert\lambda\rangle\) with \(k\)
fermionic generators having \(\sigma_{23}=+1/2\) as well as
\(\sigma_{d1}=+1/2\). (It can be shown that none of these states vanishes,
because acting on them with the adjoints of \(k\) of these fermionic
generators gives back the state \(\lvert\lambda\rangle\).) If there is a
total of \(\mathcal N\) fermionic supersymmetry generators of all types,
then there are \(\mathcal N/4\) of them with \(\sigma_{23}=+1/2\) and
\(\sigma_{d1}=+1/2\), and since these operators all anticommute the number
of states formed in this way with \(J_{23}=\lambda-k/2\) will be given by
the binomial coefficient
\begin{equation}
  \binom{\mathcal N/4}{k}\,,
  \tag{32.2.1}\label{eq:32.2.1}
\end{equation}
which when summed from \(k=0\) to the maximum value \(k=\mathcal N/4\)
gives a total of \(2^{\mathcal N/4}\) components. The minimum eigenvalue
of \(J_{23}\) obtained in this way is \(\lambda-\mathcal N/8\), reached
by multiplying the state \(\lvert\lambda\rangle\) by \(k=\mathcal N/4\)
supersymmetry generators. Taking \(\lambda=j\), we see that to avoid
having eigenvalues of \(J_{23}\) greater than \(+2\) or less than \(-2\),
we must have \(j\leq2\) and \(j-\mathcal N/8\geq-2\), which requires a
total number of fermionic generators \(\mathcal N\) no greater than 32.

Further, for \(\mathcal N=32\) supersymmetry generators there can be at
most a single supermultiplet of massless particles formed in this way by
acting on the state \(\lvert2\rangle\) with products of supersymmetry
generators having \(\sigma_{23}=+1/2\) and \(\sigma_{d1}=+1/2\). These
states have eigenvalues for any generator of the little group running in
steps of \(1/2\) from \(-2\) to \(+2\). It is only for
\(\mathcal N<32\) that there can be `matter' supermultiplets,
supermultiplets that do not contain the graviton.

A single fundamental spinor representation in \(2n\) or \(2n+1\)
dimensions has \(2^n\) components, so in order to have no more than 32
fermionic generators we must have \(n\leq5\). The spacetime dimensionality
can thus be no larger than \(d=11\) and in this case must have \(N=1\).
Supergravity in 11 dimensions is of special interest, because it may be
the `low-energy' limit of a fundamental theory known as
\(M\) theory,~\hyperref[ch32-ref-4]{[4]} which is also believed to yield
various string theories in other limits. We will now work out the spin
content of \(N=1\) supersymmetry in \(d=11\) dimensions in detail, as an
example of how this can be done by enlightened counting.

We can construct all the states of the massless multiplet for \(d=11\)
by acting on an eigenstate \(\lvert2\rangle\) of \(J_{23}\) having
eigenvalue 2 with products of \(k=0,1,\ldots,8\) supersymmetry generators
having \(\sigma_{23}=+1/2\) and \(\sigma_{2n-1,2n}=+1/2\). According to
Eq. \eqref{eq:32.2.1}, we obtain one state each with \(J_{23}=\pm2\),
eight states each with \(J_{23}=\pm3/2\), twenty-eight states each with
\(J_{23}=\pm1\), fifty-six states each with \(J_{23}=\pm1/2\), and
seventy states with \(J_{23}=0\).

For \(d=11\) the spin 2 graviton representation of the little group
\(O(9)\) is a symmetric traceless tensor with
\(9\times10/2-1=44\) independent components: there is one
\(2\mathbin{\pm}i3,2\mathbin{\pm}i3\) component with
\(J_{23}=\pm2\); seven \(2\mathbin{\pm}i3,k\) components with
\(J_{23}=\pm1\); and twenty-eight \(k,\ell\) components with
\(J_{23}=0\). (Here \(k\) and \(\ell\) run over the seven values
\(4,5,\ldots,10\). We do not count the
\(2+i3,2-i3\) component because it is related to the \(k,\ell\)
components by the tracelessness condition in this representation.)

There is also a single spin \(3/2\) gravitino representation. This
consists of a spinor \(\psi_i\) with an extra nine-vector index \(i\),
subject to an irreducibility condition \(\sum_i\gamma_i\psi_i=0\) which
excludes spin \(1/2\) components, and so has
\(9\times16-16=128\) independent components.

By subtracting the number of components with each value of \(J_{23}\)
contained in the graviton and gravitino states from those formed acting
on \(\lvert2\rangle\) with supersymmetry generators, we see that we need
one or more additional states having a total of \(28-7=21\) components
with \(J_{23}=\pm1\) and \(70-28=42\) components with \(J_{23}=0\).
The only representations of orthogonal groups that have no eigenvalues
of \(J_{ij}\)s other than \(\pm1\) and 0 are the antisymmetric tensors. An
antisymmetric tensor \(T_{i_1\cdots i_p}\) of rank \(p\) in nine
dimensions has
\[
  \begin{aligned}
    \binom{7}{p}
      &\ \text{components }T_{k_1\cdots k_p}
      &&\text{with }J_{23}=0\,,\\
    \binom{7}{p-1}
      &\ \text{components }T_{2\pm i3\,k_2\cdots k_p}
      &&\text{with }J_{23}=\pm1\,,\\
    \binom{7}{p-2}
      &\ \text{components }T_{2+i3\,2-i3\,k_3\cdots k_p}
      &&\text{with }J_{23}=0\,,
  \end{aligned}
\]
where \(k_1,\ldots,k_p\) run over the seven values \(4,5,\ldots,10\).
For \(O(9)\) the only independent antisymmetric tensors are of rank
\(p=0,1,2,3,\) and 4. The antisymmetric tensor of rank 4 has 35
components with \(J_{23}=\pm1\), which is more than needed, so it must be
excluded. Any combination of \(p=0\), \(p=1\), and \(p=2\) tensors with
the twenty-one needed components with \(J_{23}=\pm1\) would have too many
components with \(J_{23}=0\) (147 for 21 one-forms and no two-forms; 120
for 14 one-forms and 1 two-forms; 53 for 7 one-forms and 2 two-forms; and
66 for no one-forms and 3 two-forms), so we must include at least one
three-form. The antisymmetric tensor with rank \(p=3\) has just 21
components with \(J_{23}=\pm1\) and 42 components with \(J_{23}=0\),
which is just what is needed. We conclude that \emph{the unique massless
particle multiplet for \(N=1\) supersymmetry in \(d=11\) contains a
graviton, a gravitino, and a particle whose states transform under the
little group as a single antisymmetric tensor of rank 3}.

There is a richer variety of possibilities for \(d=10\). Here there are
two ways to have \(\mathcal N=32\) generators: the fermionic generators
can comprise two 16-component Weyl spinors of the same chirality, with an
automorphism group \(O(2)\), or two of opposite chirality, with no
automorphism group. For \(d=10\) it is also possible to have a single
Weyl fermionic generator, with just \(\mathcal N=16\) independent
components. These three possibilities play an important role in modern
superstring theories---they represent the massless particle spectrum of
three kinds of superstring: type IIA for 16 generators of each chirality;
type IIB for 32 generators of the same chirality; and the heterotic
superstring for 16 generators of just one chirality.

The type IIA case of \(d=10\) and opposite chirality is just like the
case of \(d=11\), except that the irreducible representations of the
little group \(O(9)\) for \(d=11\) break up into separate irreducible
representations of the little group \(O(8)\) for \(d=10\). Thus the
\(O(9)\) graviton multiplet is decomposed into an \(O(8)\) graviton with
\(8\times9/2-1=35\) components, an \(O(8)\) vector with 8 components,
and a scalar with 1 component; the \(O(9)\) gravitino multiplet is
decomposed into \(O(8)\) gravitinos of each chirality with
\((16\times8-16)/2=56\) components each and \(O(8)\) spinors of each
chirality with 8 components each; and the \(O(9)\) three-form is
decomposed into an \(O(8)\) three-form with 56 components and an \(O(8)\)
two-form with 28 components.

In the type IIB case of \(d=10\) with \(N_+=2\) and \(N_-=0\) we must
classify states according to the representation of the little group
\(O(8)\) and a quantum number \(q\) that labels the representations of
the automorphism group \(O(2)\), under which the supersymmetry generators
transform as a 2-vector. Since there is only one graviton, it must have
\(q=0\). Acting on these states with a supersymmetry generator gives
\emph{two} gravitinos with \(q=\pm1\) and 56 components each; acting with
another supersymmetry generator gives 2 two-form tensors with \(q=\pm2\)
and 28 components each; acting with another supersymmetry generator gives
2 Weyl spinors with \(q=\pm3\) and 8 components each; and acting with
another supersymmetry generator gives 2 scalars with \(q=\pm4\) together
with a self-dual four-form with \(q=0\) and 35 components.

In the heterotic case of \(d=10\) with a single Weyl fermionic generator
there are just \(\mathcal N=16\) independent components. In this case
there is a graviton supermultiplet consisting of a graviton transforming
under \(O(8)\) as a symmetric traceless tensor with 35 independent
components; a gravitino with 56 independent components; an \(O(8)\)
two-form with 28 independent components; a Weyl spinor with 8 components;
and a scalar. (This graviton supermultiplet is constructed by acting with
supersymmetry generators on one state \(\lvert2\rangle\), six states
\(\lvert1\rangle\), and one state \(\lvert0\rangle\), giving altogether
\(8\times2^4=128=35+56+28+8+1\) components.) Here we also have the
possibility of gauge supermultiplets that contain no particles having
values greater than 1 or less than \(-1\) for the eigenvalue of any
\(J_{ij}\). These gauge supermultiplets are formed by acting with
supersymmetry generators on a state \(\lvert1\rangle\), and contain one
gauge particle belonging to the vector representation of \(O(8)\), with
8 components, and one particle transforming as a fundamental Weyl spinor
of \(O(8)\), also with 8 components.
